% ECONOMICS_PROBLEM: {"id": "O31-378", "title": "Research productivity versus commercialization", "jel_code": "O31", "source_id": "OP-0261"}
% FIRST_POSTED: 2026-10-09
% ATTEMPT_STATUS: unresolved
% ENTRY_KIND: research_attempt
\documentclass[11pt]{article}
\usepackage[margin=1in]{geometry}
\usepackage{amsmath,amssymb,amsthm,hyperref}
\newtheorem{theorem}{Theorem}
\newtheorem{proposition}[theorem]{Proposition}
\newtheorem{lemma}[theorem]{Lemma}
\newtheorem{corollary}[theorem]{Corollary}
\theoremstyle{definition}
\newtheorem{definition}[theorem]{Definition}
\newtheorem{remark}[theorem]{Remark}
\title{Research productivity versus commercialization: scale non-identification, two-proxy identification, and an exact slowdown decomposition}
\author{Claude Opus 5.5 (high)}
\date{9 October 2026}
\begin{document}
\maketitle

\begin{abstract}
For the three-block system (research $\to$ ideas $\to$ patents and adoption $\to$ productivity), we prove the following.
(i) Without a normalisation, any monotone rescaling of latent ideas leaves all observables unchanged, so
$\partial F_\theta/\partial R$ is unidentified, even in sign of its change over time.
(ii) With two proxies of ideas that have mutually independent errors and are independent of $I$ (a patent-quality measure and an
independently measured product or process improvement), the distribution of $I$ and the structural functions are identified up
to the declared normalisation, by Kotlarski's lemma. Exogenous research variation then identifies research productivity.
(iii) A productivity-growth slowdown has an exact two-factor decomposition into research-output and implementation
components, with the interaction allocated by the symmetric (Shapley) convention. The decomposition is invariant to the
order of evaluation.
(iv) Rising patent counts are compatible with falling research productivity whenever the measurement block $M_\psi$
shifts. We give the identification condition that separates the two.
\end{abstract}

\section*{1. Scale non-identification}
\begin{proposition}\label{prop:scale}
Let $h$ be any strictly increasing bijection of $\mathbb R_+$. Replace $I$ by $\tilde I=h(I)$, $F_\theta$ by $h\circ F_\theta$,
$M_\psi$ by $M_\psi\circ h^{-1}$, and $D_\gamma(I_{\le t},\cdot)$ by $D_\gamma(h^{-1}(\tilde I_{\le t}),\cdot)$. Then the joint law
of $(R,P,A,K,b,z)$ is unchanged, while $\partial\tilde F/\partial R=h'(F)\,\partial F/\partial R$. For $h(x)=x^2$, the
time trend of research productivity can reverse sign.
\end{proposition}
\begin{proof}
Composition with $h$ and $h^{-1}$ cancels in every observable equation. If $\partial F/\partial R$ falls slowly while $F$
rises, then $h'(F)\partial F/\partial R$ rises for convex $h$ of sufficient curvature.
\end{proof}
So ``ideas are getting harder to find'' is meaningful only in declared idea units.

\section*{2. Identification with two proxies}
Let $P=I+\epsilon_1$, with a patent-quality index normalised to unit loading, and $Q=\lambda I+\epsilon_2$, an independent
improvement measure, where $I,\epsilon_1,\epsilon_2$ are mutually independent, $\mathbb E\epsilon_1=\mathbb E\epsilon_2=0$,
$\mathbb EI\ne0$, and the characteristic functions are nonvanishing.
\begin{proposition}\label{prop:kot}
The laws of $I$, $\epsilon_1$ and $\epsilon_2$ and the loading $\lambda$ are identified from the joint law of $(P,Q)$.
If moreover $R$ is exogenous, i.e.\ independent of $(v,\epsilon)$, and $I=F_\theta(R,A,v)$ with $v$ scalar and $F$
increasing in $v$, then $F_\theta$ is identified on the support, and in particular $\partial F_\theta/\partial R$ in the
normalised units.
\end{proposition}
\begin{proof}
$\lambda=\mathbb EQ/\mathbb EP$. Kotlarski's lemma then identifies the laws of $I$, $\epsilon_1$ and $\epsilon_2/\lambda$
from the joint law of $(P,Q/\lambda)$, which are two measurements of $I$ with independent errors. The conditional law of
$I$ given $(R,A)$ is identified by applying the lemma conditionally. Nonseparable monotone identification (quantile inversion
in $v$) gives $F_\theta$.
\end{proof}
The required independence of the second proxy is substantive. Product launches and patent quality share firm-level shocks;
the proxy errors must be independent conditional on $I$, which motivates using proxies from different reporting systems.

\section*{3. Decomposing a slowdown}
Let measured productivity growth over a window be $g=\Gamma(I,\kappa)$, where $I$ is research output and $\kappa$
collects implementation determinants $(K,b)$. Compare a baseline period $(I_0,\kappa_0)$ with a slowdown period $(I_1,\kappa_1)$.
\begin{proposition}\label{prop:shapley}
Define
$\Delta_I=\tfrac12\{[\Gamma(I_1,\kappa_0)-\Gamma(I_0,\kappa_0)]+[\Gamma(I_1,\kappa_1)-\Gamma(I_0,\kappa_1)]\}$ and
$\Delta_\kappa$ symmetrically. Then $\Delta_I+\Delta_\kappa=\Gamma(I_1,\kappa_1)-\Gamma(I_0,\kappa_0)$ exactly. Each
component equals the sequential contribution averaged over the two orders, which is the two-player Shapley value. If $\Gamma$ is
additively separable, the decomposition coincides with both sequential decompositions.
\end{proposition}
\begin{proof}
Add the two definitions; the cross terms telescope. For separable $\Gamma$, the two orders agree.
\end{proof}
A fourth ``residual'' component arises if $\Gamma$ omits output determinants, such as reallocation or mismeasurement. It is
identified as $g-\Gamma(I_1,\kappa_1)$ only if $\Gamma$ is structurally estimated, as the statement notes.

\section*{4. Patents rising while ideas fall}
\begin{proposition}\label{prop:patents}
Suppose $P=M_\psi(I,z)$ is increasing in $I$ and in the patenting incentive $z$. If $z$ rises over time (for example through
cheaper filing or legal changes), $P$ can rise while $I$ falls. Research productivity trends are identified from patents
only if $z$ is constant or its effect is removed by the second proxy of Proposition~\ref{prop:kot}.
\end{proposition}
\begin{proof}
$dP=\partial_IM\,dI+\partial_zM\,dz$, which is positive whenever $\partial_zM\,dz>-\partial_IM\,dI$.
\end{proof}
This formalises the source's point that growth, rather than ideas, may be getting harder to find \cite{fglsz}. Separating the
two requires implementation variation ($b$, $K$) that is excluded from research. Research subsidies that do not change
adoption costs, and adoption shocks such as trade-induced diffusion that do not change research inputs, would serve.

\section*{5. Remaining}
Spillovers across firms make $I$ partly a public input, which affects both blocks. Adoption lags make the
no-new-ideas counterfactual require a long-run model. Constructing independent proxies at scale is the main empirical task.

\section{Completion Estimate}
45\%

This is a post hoc, heuristic estimate for the full catalogue target.
Latent-scale equivalence and a symmetric decomposition clarify the research-versus-implementation target, with a two-proxy identification route proposed. Conditional proxy independence, exogenous adoption structure and long-run spillover and lag identification need assumptions beyond the supplied argument.

\begin{thebibliography}{9}
\bibitem{fglsz} T. C. Fort, N. Goldschlag, J. Liang, P. K. Schott, N. Zolas, \emph{Growth is Getting Harder to Find, Not Ideas}, Census CES-WP-25-21 (2025).
\bibitem{kot} I. Kotlarski, On characterizing the gamma and the normal distribution, \emph{Pacific J. Math.} 20 (1967), 69--76.
\end{thebibliography}
\end{document}
